% Source: https://artofproblemsolving.com/community/c3h1373600_confused_about_a_solution
% Stable identity: aops-topic-1373600-q01
From $(-3,0)$ to $(0,k+3)$ the run is $3$ and the rise is $k+3$, so the slope is $\dfrac{k+3}{3}$. Setting it equal to $k$:
\[
\dfrac{k+3}{3}=k\iff k+3=3k\iff k=\boxed{\dfrac{3}{2}}.
\]

\begin{asy}
usepackage("amsmath");
size(5.5cm);
defaultpen(fontsize(10pt));
import graph;
real k=1.5;
draw((-4.5,0)--(2,0),Arrow(6)); draw((0,-1)--(0,6.5),Arrow(6));
label("$x$",(2,0),E); label("$y$",(0,6.5),N);
draw((-4,-1.5)--(1,6),blue+linewidth(0.9));
pair P=(-3,0), Q=(0,4.5);
dot(P); dot(Q);
label("$(-3,0)$",P,SE); label("$(0,k+3)$",Q,E);
draw(P--(0,0),red+dashed); draw((0,0)--Q,red+dashed);
label("run $3$",(-1.5,0),N,red); label("rise $k+3$",(0,2.25),E,red);
label("\mbox{}",(min(currentpicture,user=true)+max(currentpicture,user=true))/2,black); // keeps the text after the figure black when LaTeX typesets the labels inline
\end{asy}

Check: with $k=\dfrac{3}{2}$ the line passes through $(-3,0)$ and $\left(0,\dfrac{9}{2}\right)$, and its slope is $\dfrac{9}{2}\div3=\dfrac{3}{2}=k$.

So the student's answer $\dfrac{3}{2}$ is right. An expression such as $\dfrac{k+3}{3}$ (not $\dfrac{k+3}{2}$) is only the slope written in terms of $k$, an intermediate step, not the value of $k$.
